\section{Data, Information Boundaries, and Systems}

\subsection{Primary Datasets and Lineage}

FinQA pairs financial-report questions with text, tables, and executable
reasoning programs \citep{chen2021finqa}. TAT-QA similarly combines tabular and
textual financial evidence with arithmetic and counting questions
\citep{zhu2021tatqa}. Preparation consumes pinned local source snapshots; network
acquisition is a separate explicit operation. The adapter verifies each snapshot
hash, safely executes the annotated derivation, locates every operand in public
evidence, and rejects unsupported operations, ambiguous scales, missing support,
duplicate documents, or answer-program disagreements.

At most one question is selected from each source document or context. A
\texttt{headroom} case requires at least two derivation operations and at least
one required item ranked 3--12 by the frozen baseline retriever. An
\texttt{easy\_control} case uses one operation and places all required evidence
in the first two results. Splits are deterministic and document-disjoint: 100
development cases, 60 operational-pilot cases, up to 1,000 balanced hard main
cases, and 100 balanced easy reserves. Quota or eligibility shortfalls abort
preparation rather than silently changing the cohort.

Public cases contain only the question, evidence corpus, outcome-independent
stratum, document identity, and descriptive metadata. Hidden labels contain the
typed answer, executable gold derivation, support IDs, and source lineage. The
two stores join only after generation. Numeric scoring parses the strict
candidate schema; it never extracts the first or last number from prose. It uses
decimal arithmetic, explicit unit and scale normalization, entity and period
compatibility, and frozen absolute and relative tolerances.

\subsection{Systems and Mechanisms}

Every system uses the same core instructions, evidence serialization, strict
candidate schema, exact model, and deterministic retriever. The manipulated
factor is inference structure.

\begin{table}[H]
\centering
\small
\caption{Frozen system mechanisms. Unverified search is exploratory only.}
\label{tab:systems}
\begin{tabularx}{\linewidth}{@{}lY@{}}
\toprule
System & Frozen mechanism \\
\midrule
Monolith & Direct question query, tier-limited retrieval, and exactly one answer
call capped at 256 output tokens. It never checks or revises. \\
Verified search & Mandatory planner capped at 128 tokens; deterministic query
union; sequential candidates capped at 256; label-blind checking; first passing
candidate; at most one middle/high repair using checker findings only. \\
Unverified search & Same planner, retrieval, and candidate opportunity as
verified search, but no checker feedback. It generates every affordable
candidate and selects normalized plurality with stable candidate-order ties. \\
\bottomrule
\end{tabularx}
\end{table}

The checker verifies citation existence, operand provenance, safe arithmetic,
candidate-expression agreement, units, scale, entity, period, and division
safety. It cannot see correctness or a gold value. Every cell records planned
and actual queries, query hashes, retrieval IDs before and after truncation,
candidates, checks, repair, accepted index, answer changes, authoritative usage,
realized tokens, and a frozen exit reason. Invalid output, refusal,
architecture-caused tool failure, budget exhaustion, and abstention score
incorrect under intention-to-treat. Infrastructure failures are unscored and
trigger whole matched-block handling.

\subsection{Action-Backed Resource Intervention}

The hard budget is the sum of authoritative prompt and completion usage over all
calls in a cell:
\begin{equation}
T_{isb}=\sum_k(P_{isk}+C_{isk})\le B_b.
\label{eq:budget}
\end{equation}
Before each call, the ledger obtains an exact chat-tokenizer prompt count and
reserves the maximum output. It refuses an unaffordable call, releases unused
reservation after the authoritative usage arrives, and treats prompt mismatch,
output-reservation overrun, or hard-cap overrun as a protocol violation.

\begin{table}[H]
\centering
\small
\caption{Action-backed budget tiers.}
\label{tab:budgets}
\begin{tabular}{lrrrrr}
\toprule
Tier & Token cap & Retrieval & Queries & Candidates & Repairs \\
\midrule
Low & 4,096 & 2 & 1 & 1 & 0 \\
Middle & 12,288 & 6 & 2 & 2 & 1 \\
High & 32,768 & 12 & 4 & 4 & 1 \\
\bottomrule
\end{tabular}
\end{table}

Outcome-free development may advance a ceiling only through 8,192, 16,384,
24,576, 32,768, 49,152, and 65,536 when more than 1\% of development cases
cannot fit mandatory prompts and actions. Accuracy and architecture differences
are unavailable during calibration. The selected ceilings freeze before pilot.
